\documentclass[12pt, a4paper]{article}
\usepackage[utf8]{inputenc}
\usepackage[spanish]{babel}
\usepackage{amsmath}
\usepackage{amssymb}
\usepackage{geometry}
\usepackage{xcolor}
\usepackage{titlesec}
\usepackage{booktabs}

\geometry{top=2.5cm, bottom=2.5cm, left=2.5cm, right=2.5cm}

% Colores personalizados
\definecolor{navyblue}{RGB}{0, 51, 102}

% Estilo de los títulos
\titleformat{\section}{\color{navyblue}\normalfont\Large\bfseries}{\thesection}{1em}{}
\titleformat{\subsection}{\color{navyblue}\normalfont\large\bfseries}{\thesubsection}{1em}{}

\title{\textbf{Fórmulas de las Prestaciones y Rendimiento de la CPU}}
\author{Eric Guacache}
\date{}

\begin{document}
	
	\maketitle
	
	\section{Fórmulas para calcular las Prestaciones de un Computador}
	
	Para calcular la Prestación de un computador $X$
	
	\begin{equation}
		\text{Prestaciones}_X = \frac{1}{\text{Tiempo de ejecución}_X}
	\end{equation}
	
	Al comparar dos máquinas distintas, $X$ e $Y$, si las prestaciones de $X$ son mayores que las de $Y$, se cumple la siguiente relación:
	\begin{equation}
		\text{Prestaciones}_X > \text{Prestaciones}_Y
	\end{equation}
	Sustituyendo la definición fundamental, obtenemos:
	\begin{equation}
		\frac{1}{\text{Tiempo de ejecución}_X} > \frac{1}{\text{Tiempo de ejecución}_Y}
	\end{equation}
	Lo que implica necesariamente que el tiempo de ejecución de la máquina de menores prestaciones es mayor:
	\begin{equation}
		\text{Tiempo de ejecución}_Y > \text{Tiempo de ejecución}_X
	\end{equation}
	Esto significa que el tiempo de ejecución de $Y$ es mayor que el de $X$, confirmando que $X$ es más rápido que $Y$.
	
	\subsection{Comparación de Prestaciones de 2 Computadores}
	Para expresar numéricamente cuánto más rápida es una máquina en comparación con otra, utilizamos la frase ``$X$ es $n$ veces más rápida que $Y$'', lo cual se define matemáticamente como:
	\begin{equation}
		\frac{\text{Prestaciones}_X}{\text{Prestaciones}_Y} = n
	\end{equation}
	Haciendo el cambio a términos de tiempo de ejecución, la relación se invierte:
	\begin{equation}
		\frac{\text{Prestaciones}_X}{\text{Prestaciones}_Y} = \frac{\text{Tiempo de ejecución}_Y}{\text{Tiempo de ejecución}_X} = n
	\end{equation}
	
	\section{Ecuaciones para calcular el tiempo de ejecución de CPU}
	
	Para calcular el tiempo de ejecución de CPU se usan las siguientes fórmulas:
	
	\begin{equation}
		\text{Tiempo de ejecución de CPU} = \text{Ciclos de reloj de la CPU} \times \text{Tiempo del ciclo de reloj}
	\end{equation}
	
	Alternativamente, dado que la frecuencia de reloj ($f$) es la inversa del tiempo de ciclo ($T = 1/f$), la ecuación se puede reescribir de la siguiente manera:
	
	\begin{equation}
		\text{Tiempo de ejecución de CPU} = \frac{\text{Ciclos de reloj de la CPU}}{\text{Frecuencia de reloj}}
	\end{equation}
	
	\section{La Ecuación alternativa para el tiempo de ejecución}
	
	También puedes usar esta ecuación alternativa para el tiempo de ejecución usando el \textbf{número de instrucciones} (instrucciones ejecutadas por el programa), el \textbf{CPI} (Ciclos Por Instrucción) y el tiempo de ciclo:
	
	\begin{equation}
		\text{Tiempo de ejecución} = \text{Número de instrucciones} \times \text{CPI} \times \text{Tiempo de ciclo}
	\end{equation}
	
	O bien, utilizando nuevamente el inverso del tiempo de ciclo (la frecuencia de reloj):
	\begin{equation}
		\text{Tiempo de ejecución} = \frac{\text{Número de instrucciones} \times \text{CPI}}{\text{Frecuencia de reloj}}
	\end{equation}
	
	\section{Cosas que tienes que recordar}
	
	\begin{table}[h!]
		\centering
		% Cambiamos lll por columnas p{} para permitir que el texto salte de línea y no desborde la página
		\begin{tabular}{l | p{7cm} | p{4.5cm}}
			\toprule
			\textbf{Métrica} & \textbf{Definición} & \textbf{Unidad Típica} \\
			\midrule
			Tiempo de ejecución & Tiempo total dedicado a una tarea & Segundos (s), Milisegundos (ms) \\
			\midrule
			Número de instrucciones & Cantidad de instrucciones ejecutadas & Instrucciones \\
			\midrule
			CPI & Ciclos de reloj promedio por instrucción & Ciclos / Instrucción \\
			\midrule
			Frecuencia de reloj & Velocidad de oscilación del reloj del sistema & Hertz (Hz), Gigahertz (GHz) \\
			\bottomrule
		\end{tabular}
	\end{table}
	
\end{document}

